\section{Sharpness, and some structural remarks}
\label{sec:sharp}

The finite theorem uses the elementary function $g_\alpha$ of \cref{lem:scalar}. This section shows that the same function is the exact limiting moment exponent, so the scalar step loses no exponential rate at large width. It then records some structural consequences of the gauge identity.

\subsection{The exact scalar high-moment exponent}

For $\theta>0$ and $0\le u\le1$ define
\begin{equation}\label{eq:Phi}
\Phi_{\alpha,\theta}(u)=H(u)-\log2+\theta\log(2\alpha)+\Bigl(\frac u2+\theta\Bigr)\log\Bigl(\frac u2+\theta\Bigr)-\frac u2\log\frac u2-\theta,
\end{equation}
continuous on $[0,1]$ because $x\log x$ extends continuously to zero.

For a chi-square variable with $m\ge1$ degrees of freedom and $r>0$, the substitution $x=2y$ in $\int_0^\infty x^r\frac{x^{m/2-1}e^{-x/2}}{2^{m/2}\Gamma(m/2)}\dd x$ gives
\begin{equation}\label{eq:chi-moment}
\E[(\chi^2_m)^r]=2^r\,\frac{\Gamma(m/2+r)}{\Gamma(m/2)} .
\end{equation}
Conditioning on $M_n$ in \cref{eq:mixture} with $r=\theta n$ (the term $m=0$ contributes zero),
\begin{equation}\label{eq:exact-sum}
\E[(X_n^{(\alpha)})^{\theta n}]=2^{-n}\sum_{m=1}^n\binom nm\Bigl(\frac{2\alpha}n\Bigr)^{\theta n}\frac{\Gamma(m/2+\theta n)}{\Gamma(m/2)} .
\end{equation}
Two forms of Stirling's formula are used: for integers $k\ge1$, $\log k!=(k+\tfrac12)\log k-k+O(1)$, and for real $x\ge\tfrac12$, $\log\Gamma(x)=(x-\tfrac12)\log x-x+O(1)$ uniformly in $x\ge\tfrac12$ (the function $\log\Gamma(x)-(x-\tfrac12)\log x+x$ is continuous on $[\tfrac12,\infty)$ and tends to $\tfrac12\log2\pi$). With $u=m/n$ they give, uniformly for $1\le m\le n$,
\begin{gather*}
\log\binom nm=nH(u)+O(\log n),\\
\log\Gamma\Bigl(n\bigl(\tfrac u2+\theta\bigr)\Bigr)-\log\Gamma\Bigl(\frac{nu}2\Bigr)=\theta n\log n+n\Bigl[\bigl(\tfrac u2+\theta\bigr)\log\bigl(\tfrac u2+\theta\bigr)-\tfrac u2\log\tfrac u2-\theta\Bigr]+O_\theta(\log n).
\end{gather*}
The term $-\theta n\log n$ from $(2\alpha/n)^{\theta n}$ cancels the $+\theta n\log n$ of the Gamma ratio, and the $m$-th summand $a_{n,m}$ of \cref{eq:exact-sum} satisfies $\frac1n\log a_{n,m}=\Phi_{\alpha,\theta}(m/n)+O_{\alpha,\theta}(\log n/n)$ uniformly. A sum of $n$ nonnegative terms lies between its largest term and $n$ times it, so after taking $\frac1n\log$ the sum and the maximum differ by at most $\log n/n$; and the grids $\{1/n,\dots,1\}$ become dense, so the grid maximum of the continuous $\Phi_{\alpha,\theta}$ converges to its maximum.

\begin{proposition}[exact exponent]\label{prop:exact}
For every $\alpha>0$ and $\theta>0$,
\[
\lim_{n\to\infty}\frac1n\log\E[(X_n^{(\alpha)})^{\theta n}]
=\max_{0\le u\le1}\Phi_{\alpha,\theta}(u)=g_\alpha(\theta),
\]
where $g_\alpha$ is the finite-width bound in \cref{eq:scalar-def}. The unique maximizing fraction is $u_\theta=1/(1+z_\theta)$.
\end{proposition}

\begin{proof}
The limit and its variational formula follow from the exact Gamma sum and uniform estimates above. For $0<u<1$,
\[
\partial_u\Phi=\log\frac{1-u}{u}+\frac12\log\frac{u/2+\theta}{u/2},\qquad
\partial_{uu}\Phi=-\frac1{u(1-u)}+\frac1{4(u/2+\theta)}-\frac1{2u}<0.
\]
The first derivative tends to $+\infty$ at zero and $-\infty$ at one, so there is one interior maximizer. Substitute $u=1/(1+z)$. Its stationary equation reduces to $\theta=(1-z)/(2z^2)$, giving $z=z_\theta$ and
\[
\frac{u_\theta}2+\theta=\frac1{2z_\theta^2(1+z_\theta)}.
\]
Substituting these identities into \cref{eq:Phi} and collecting logarithms gives exactly \cref{eq:scalar-def}. This identifies the limiting exponent with the function that already bounds every finite-width moment in \cref{lem:scalar}.
\end{proof}

In particular $u_\theta=\tfrac12+\tfrac12\theta+O(\theta^2)$ and \cref{eq:scalar-expansion} gives
\begin{equation}\label{eq:g-expansion}
g_{2e^{-\delta}}(\theta)=-\delta\theta+\frac52\theta^2+O(\theta^3).
\end{equation}
The scalar exponent is exact, but the geometric count and the sphere net in the master inequality still introduce slack. The constant $\sqrt{10}$ remains a constant in a certified lower bound for the network threshold, not a matching asymptotic description of that threshold.

\subsection{The Cover factor at fixed depth}

\begin{lemma}\label{lem:C-rate}
For every fixed $L\ge2$, $\lim_{n\to\infty}\frac1n\log C_{n,L}=LH(1/L)$.
\end{lemma}

\begin{proof}
For $L>2$ the ratio of consecutive terms $\binom{nL-1}{j+1}/\binom{nL-1}{j}=(nL-1-j)/(j+1)$ exceeds $1$ for $0\le j\le n-2$, because $nL-1-j\ge nL-n+1>n-1\ge j+1$; so the last term is the largest and $\binom{nL-1}{n-1}\le C_{n,L}/2\le n\binom{nL-1}{n-1}$. Stirling's formula gives $\frac1n\log\binom{nL-1}{n-1}\to LH(1/L)$, and the factors $2$ and $n$ vanish after $\frac1n\log$. For $L=2$, $C_{n,2}=2^{2n-1}$ and $\frac1n\log C_{n,2}\to2\log2=2H(\tfrac12)$.
\end{proof}

So the entropy bound of \cref{lem:C-entropy} is exact at the exponential scale, and the finite theorem loses nothing there at large width. At small width it does lose something; \cref{sec:numerics} shows how much.

\subsection{The all-active pattern, exactly}

Take $D_1=\dots=D_L=I_n$. Then $H^D_\ell=Q_\ell:=W_\ell\cdots W_1$ and $A_D$ stacks $Q_1,\dots,Q_L$. Every Gaussian square matrix is invertible almost surely, since its determinant is a nonzero polynomial in continuously distributed entries and the zero set of a nonzero polynomial has Lebesgue measure zero (\cref{lem:poly}). On that event $(W_1,\dots,W_L)\mapsto(Q_1,\dots,Q_L)$ is a smooth bijection with inverse $W_1=Q_1$, $W_\ell=Q_\ell Q_{\ell-1}^{-1}$, so the tuple $(Q_\ell)$ has a joint density on $(GL_n)^L$. Choose any $n$ of the $nL$ rows of the stacked matrix; their determinant is a polynomial in the entries of the $Q_\ell$, not identically zero (assign the chosen rows the standard basis vectors and complete each $Q_\ell$ to an invertible matrix), hence nonzero almost surely; finitely many choices, so all $n\times n$ minors are nonzero almost surely. That is central general position for the coordinate hyperplanes restricted to $\im A_D$, and the region count is then exact: $r(\im A_D)=C_{n,L}$ almost surely. \Cref{prop:orbit} with $F=1$ gives
\begin{equation}\label{eq:all-active}
\PP(R_D\ne\varnothing)=\frac{C_{n,L}}{2^{nL}}=2^{1-nL}\sum_{j=0}^{n-1}\binom{nL-1}j\qquad(D\text{ all-active}),
\end{equation}
an exact formula for the probability that a random deep network has an input at which every unit is active at every layer. Its logarithm is $-[nL\log2-\log C_{n,L}]=-n[L\log2-LH(1/L)]+o(n)$ by \cref{lem:C-rate}: the accumulated cost per unit width is $L\log2-O(\log L)$, not a constant per layer. For $L=1$ the cost is zero, as it must be: $C_{n,1}=2^n$, and a full-rank $W_1$ reaches every orthant. \Cref{sec:numerics} checks \cref{eq:all-active} against direct enumeration.

\begin{lemma}\label{lem:poly}
If $P:\R^d\to\R$ is a polynomial that is not identically zero, then $\{x:P(x)=0\}$ has $d$-dimensional Lebesgue measure zero.
\end{lemma}

\begin{proof}
Induction on $d$. For $d=1$ a nonzero polynomial has finitely many roots. For $d>1$ write $P=\sum_{k=0}^ma_k(x_1,\dots,x_{d-1})x_d^k$ with some $a_k$ not identically zero. By induction the set where all $a_k$ vanish has measure zero; for every other value of the first $d-1$ coordinates the polynomial in $x_d$ is nonzero and has finitely many roots, a one-dimensional null set. Fubini's theorem gives total measure zero.
\end{proof}

\subsection{Depth asymptotics at fixed width}
\label{sec:fixed-width-depth}

The same identity also controls the opposite order of limits: keep the width fixed and let the depth grow. Write $R_{n,L}=\#\{D:R_D\ne\varnothing\}$ for the strict-region count and $S_{n,L}=\{F^{(\alpha)}_{n,L}\not\equiv0\}$ for global survival. A strictly negative final preactivation can leave a nonempty strict region whose output is zero, so these are different objects.

\begin{corollary}[fixed-width depth asymptotics]\label{cor:fixed-width-depth}
Fix $n\ge1$ and $\alpha>0$, with independent Gaussian layers, no biases, input domain $\R^n$ and ReLU at the output as in \cref{eq:forward}. Put $\beta_n=1-2^{-n}$. There are finite positive constants $\kappa_n$ and $\sigma_n$, independent of $\alpha$, such that
\begin{equation}\label{eq:depth-ratios}
\E R_{n,L}\sim\kappa_n\beta_n^L,
\qquad \PP(S_{n,L})\sim\sigma_n\beta_n^L
\qquad (L\to\infty).
\end{equation}
Both ratios on the left after division by $\beta_n^L$ are nondecreasing. Moreover $\kappa_n\ge n\beta_n^{-2}$ and $\beta_n^{-1}\le\sigma_n\le\kappa_n$. In particular,
\begin{equation}\label{eq:depth-rate}
\lim_{L\to\infty}\frac{\log\E R_{n,L}}L
=\lim_{L\to\infty}\frac{\log\PP(S_{n,L})}L
=\log(1-2^{-n}).
\end{equation}
\end{corollary}

\begin{proof}
Write $q=2^n$, $m=q-1$, $a=q-n-1$ and $\beta=m/q$. For every fixed sequence of nonzero intermediate masks, every row of every formal preactivation matrix is nonzero almost surely. Indeed, conditional on a nonzero retained product $B$, a fresh Gaussian row $w$ satisfies $w^TB\ne0$ almost surely: its zero set is a proper linear subspace. A nonzero mask retains at least one such row. Induction, followed by a finite union over rows and masks, proves the assertion at each depth. These are formal masks; no conditioning on realizability is involved.

Let $A_k$ be the scalar sum of $\E r(\im A_D)$ for the prefix stack $H^D_1,\ldots,H^D_k$, over all $k-1$ intermediate masks of rank at least two. The final mask is omitted because it does not enter the stack. There are $a^{k-1}$ prefixes, and the invertibility of $W_1$ gives
\[
A_1=q,\qquad 0\le A_k\le a^{k-1}C_{n,k}.
\]
At $n=1$, use $0^0=1$: only the empty prefix exists, and $A_k=0$ for $k>1$.

Partition the nonzero intermediate-mask sequences by their first rank-one mask, at layer $k<L$. The selected row of $H^D_k$ already occurs in the prefix stack. Every later row is a nonzero multiple of it almost surely, so later layers add no new hyperplane and leave the orthant count unchanged. There are $n$ choices for the selected row, $m^{L-k-1}$ choices for the later intermediate masks, and $q$ choices for the final mask. This family contributes $nq m^{L-k-1}A_k$ to the formal orthant sum. The family with no rank-one intermediate mask contributes $qA_L$; a zero intermediate mask contributes zero because the next preactivation block vanishes. Summing \cref{eq:orbit} with functional $1$ therefore gives the exact identity
\begin{equation}\label{eq:depth-prefix}
\E R_{n,L}=q^{1-L}\left[A_L+n\sum_{k=1}^{L-1}m^{L-k-1}A_k\right],
\qquad L\ge1.
\end{equation}
After division by $\beta^L$ this is
\[
\frac{\E R_{n,L}}{\beta^L}
=\frac{qA_L}{m^L}+\frac{nq}{m^2}\sum_{k=1}^{L-1}\frac{A_k}{m^{k-1}}.
\]
For $n\ge2$, $\rho=a/m<1$, and $A_k/m^{k-1}\le C_{n,k}\rho^{k-1}$. At fixed $n$, $C_{n,k}$ is a polynomial in $k$ of degree $n-1$. The first term tends to zero and the series converges, proving
\begin{equation}\label{eq:depth-constant}
\kappa_n=\frac{nq}{m^2}\sum_{k\ge1}\frac{A_k}{m^{k-1}}\in(0,\infty).
\end{equation}
The first summand gives $\kappa_n\ge nq^2/m^2=n\beta^{-2}$. The same formula gives $\kappa_1=4$ directly. Extending any prefix by one of the $a$ admissible masks and a fresh layer only appends proper hyperplanes, so it cannot remove an old open region. Thus $A_{k+1}\ge aA_k$, and subtraction in \cref{eq:depth-prefix} yields
\[
\E R_{n,L+1}-\beta\E R_{n,L}
=q^{-L}(A_{L+1}-aA_L)\ge0.
\]
This proves monotonicity of the normalized means.

For survival, couple the depths by one infinite sequence of independent matrices and put $p_L=\PP(S_{n,L})$. On $S_{n,L}$, choose the first input in a fixed enumeration of $\mathbb Q^n$ at which $F^{(\alpha)}_{n,L}$ is nonzero. Such an input exists by continuity, and the choice is measurable with respect to $W_1,\ldots,W_L$ because the enumeration is countable. Conditional on this past, the next fresh Gaussian layer preserves its nonzero output with probability exactly $\beta$. Its survival implies global survival, so $p_{L+1}\ge\beta p_L$; equality is not asserted. Positive homogeneity gives $F(0)=0$, and \cref{lem:reduction} ensures that a nonzero network has a realized strict region. Consequently $p_L\le\E R_{n,L}$. The sequence $p_L/\beta^L$ is therefore nondecreasing and bounded above by $\kappa_n$. Since $p_1=1$ almost surely, it converges to $\sigma_n\in[\beta^{-1},\kappa_n]$. Finally, a positive common rescaling of Gaussian weights changes neither strict signs nor zero-output events, so both constants are independent of $\alpha>0$.
\end{proof}

The finite bounds behind these asymptotics are also explicit:
\begin{align}
n\beta_n^{L-2}\ \le\ \E R_{n,L}\ &\le\ C_{n,L}\beta_n^{L-1}, &&L\ge2,\label{eq:depth-strict}\\
\beta_n^{L-1}\ \le\ \PP(S_{n,L})\ &\le\ \min\{1,C_{n,L}\beta_n^L\}, &&L\ge1.\label{eq:depth-survival}
\end{align}
For the mean, the lower bound is the $k=1$ term in \cref{eq:depth-prefix}; the upper bound counts the $q m^{L-1}$ masks with nonzero intermediate masks and bounds each orthant count by $C_{n,L}$. For global survival, use \cref{prop:orbit} with the invariant functional $\one_{\{J_D\ne0\}}$: a zero mask at any layer, including the final one, kills the Jacobian, leaving at most $m^L$ contributors. For the lower bound choose $s=W_1^{-1}e_1$, so $x_1(s)=e_1$, and propagate this first-layer witness through the independent remaining layers. Its survival probability is $\beta_n^{L-1}$. A deterministic nonzero input instead survives with probability exactly $\beta_n^L$.

At width one, direct sign enumeration gives $\E R_{1,L}=2^{2-L}$ and $\PP(S_{1,L})=2^{1-L}$, hence $(\kappa_1,\sigma_1)=(4,2)$. At $L=2$ there is one strict region for every nonzero scalar pair, while global survival has probability $1/2$. The constants at higher widths are left implicit, except for $\kappa_2=8$ below. These are fixed-width limits; no uniformity in a growing width or quantitative convergence rate for the survival ratio is asserted.

Earlier work already establishes death at large depth. In the zero-bias square setting, Rister and Rubin~\cite{rister2021}, Sections~3--4, give a global survival upper bound with base $1-2^{-n^2}$ and a fixed-input lower bound with base $\beta_n$. Their Section~5.3 also studies a survival ratio for finite data as width grows, under a variance assumption. Lu et al.~\cite{lu2020}, Theorem~3.5 and Appendix~C, give two-exponential bounds and a one-ray state analysis for one-dimensional input, with a final affine readout. After accounting for that depth convention, their bounds and the past-measurable witness argument above also imply a finite positive survival-ratio limit in the scalar-input setting. For any fixed finite set of $M$ nonzero inputs, the simpler bound $p_L\le M\beta_n^L$ and the same monotonicity argument suffice. The additional geometric step here bounds the normalized mean strict count and hence survival on the entire square input space $\R^n$. Neither qualitative eventual death nor ratio arguments in general are claimed as new.

\subsection{The mean strict-region count at width two}

Let $R_{2,L}$ be the number of nonempty strict regions $R_D$ in a width-two network. This counts activation patterns with every preactivation nonzero; it excludes ordinary cells on which a preactivation is identically zero.

\begin{proposition}\label{prop:width-two-mean}
For every $L\ge1$ and every $\alpha>0$,
\begin{equation}\label{eq:width-two-mean}
\E R_{2,L}=8\left[\left(\frac34\right)^L-\left(\frac14\right)^L\right].
\end{equation}
In particular, $\kappa_2=8$ in \cref{eq:depth-ratios}.
\end{proposition}

\begin{proof}
At width two, the only intermediate mask of rank at least two is $I_2$. Thus the prefix in \cref{eq:depth-prefix} stacks the unmasked products $Q_1,\ldots,Q_k$. Their $2k$ row directions are distinct almost surely by the joint-density argument for \cref{eq:all-active}. They define $2k$ lines in the input plane, so $A_k=4k$. Substituting $q=4$, $m=3$ and $n=2$ in \cref{eq:depth-prefix} gives
\[
\E R_{2,L}=4^{-L}\left[32\sum_{k=1}^{L-1}k3^{L-k-1}+16L\right]
=8\frac{3^L-1}{4^L}.
\]
The finite-sum identity follows by differentiating a geometric sum, including $L=1$ when the sum is empty. Dividing by $(3/4)^L$ gives the limit $8$.
\end{proof}

The first six means are $4,4,3.25,2.5,1.890625,1.421875$. Their decay at large depth concerns strict regions. Ordinary cells remain even when the network is identically zero, so \cref{eq:width-two-mean} must not be read as their mean count. No affine-bias version of this formula is asserted.

\subsection{Depth one}

At $L=1$ the all-active mask is realized with probability one, $J_D=W_1$ for it, and \cref{eq:reduction} gives $K_{n,1}(\alpha)=\op{W_1}$: every mask $D_1$ gives $\op{D_1W_1}\le\op{W_1}$. By the Bai--Yin theorem \cite{baiyin}, for the centered Gaussian matrix of this paper, with independent $\mathcal N(0,\alpha/n)$ entries, $\op{W_1}\to2\sqrt\alpha$ almost surely. Hence $\PP(K_{n,1}(\alpha)\le1)\to1$ for $\alpha<\tfrac14$ and $\to0$ for $\alpha>\tfrac14$, that is,
\[
\alpha_1=\tfrac14 .
\]
A depth-one network at $\alpha=1$ is expansive with probability tending to one, although $1<2$; there is no contradiction with \cref{thm:main}, whose part (i) only takes effect beyond $L_0(\alpha)$.

\subsection{Why the global upper bound never conditions on an input}

A tempting route would be to choose an input, condition on the activation masks it generates, and analyse the resulting Jacobian as if the masks were independent of the weights. The last step is false in general: the event that a mask is generated by an input carries information about the same weights that appear in the Jacobian. The gauge argument fixes a formal mask without choosing a witness input, multiplies the realizability indicator by a gauge-invariant Jacobian functional, and averages over a finite symmetry group exactly; only then does it sum over masks. The sphere net is used for tangent directions of a fixed linear map, not to search for regions, so a realizable cone of arbitrarily small angle is still counted.

\subsection{The order of limits}

The master inequality \cref{eq:master} and the finite theorem are statements about every $n$ and $L$, and one may evaluate them in any joint regime. What is proved about the thresholds, however, fixes $L$ first. If $L=L(n)$ grows rapidly, several effects change scale; for instance the probability that some layer dies at a fixed input is at most $L2^{-n}$, negligible for fixed $L$ but not for every $L(n)$. No joint-limit statement should be read into \cref{thm:main} beyond what \cref{thm:finite} gives when its budget is evaluated in that regime.
